% SPDX-License-Identifier: AGPL-3.0-or-later
% reflowtex capture, read at \begin{document} (begindocument/before), after
% the document's whole preamble: what has to see the preamble's
% definitions – a class's or package's \ref, \label, \footnote, \marginpar,
% sectioning and theorems, TikZ if the preamble loaded it – and the layout
% the preamble chose. It is read as the document's own code (LaTeX's
% top-level hook label), so what it adds to the hooks of \begin{document}
% runs after every package's.

% A web page has no columns. A two-column document is set in one, at its full
% \textwidth, so its paragraphs are as wide as the measure the viewer re-breaks
% them from (set at the column instead, they would keep the other column's
% width as a margin). The switches a document uses later – \twocolumn[...]
% for a title block, revtex's \twocolumngrid after its front matter – stay in
% one column too.
\makeatletter
\AtBeginDocument{%
  \if@twocolumn
    \global\@twocolumnfalse
    \columnwidth=\textwidth \hsize=\textwidth \linewidth=\textwidth
  \fi
  \@ifundefined{onecolumngrid}{}{%
    \onecolumngrid
    \global\let\twocolumngrid\onecolumngrid}%
  \global\let\reflowtexOrigOnecolumn\onecolumn
  \gdef\twocolumn{\@ifnextchar[{\reflowtexOneColumnWith}{\reflowtexOrigOnecolumn}}}
\long\def\reflowtexOneColumnWith[#1]{\reflowtexOrigOnecolumn #1\par}
\makeatother

% A table of contents, rebuilt for a page. A pageless document never turns a
% page, so the page numbers it would print are all 1: they are left out. And
% each entry becomes a link to its heading: the k-th line read from the .toc
% was written by the k-th \addcontentsline{toc} of the run before, so it links
% to the anchor that one places this run (reflowtex-toc-k, in the outline
% capture below). A .toc line has a fourth argument when whatever wrote it
% reads one (hyperref's anchor name, the kernel's destination since 2023), and
% it is passed on untouched exactly then.
\newcount\reflowtexTocEntryCount
\newcount\reflowtexTocAnchorCount
\makeatletter
\let\reflowtex@tocfile\@empty
\def\reflowtex@newstarttoc#1{%
  \def\reflowtex@tocfile{#1}\reflowtex@starttoc{#1}\let\reflowtex@tocfile\@empty}
\def\reflowtex@newcontentsline#1#2#3{%
  \@ifnextchar\bgroup{\reflowtex@contentsline{#1}{#2}}{\reflowtex@contentsline{#1}{#2}\reflowtex@noarg}}
\def\reflowtex@noarg{\reflowtex@noarg}   % stands for a fourth argument the line did not have
\def\reflowtex@contentsline#1#2#3{%
  \def\reflowtex@tmp{toc}%
  \ifx\reflowtex@tocfile\reflowtex@tmp
    \global\advance\reflowtexTocEntryCount by 1
    \global\advance\reflowtexLinkCount by 1
    \directlua{Serializer.note_link(\the\reflowtexLinkCount,
      "reflowtex-toc-\the\reflowtexTocEntryCount")}%
    % no group around the entry: revtex's \l@section hands the widest
    % number it measured to \aftergroup, which a group of ours would eat
    \edef\reflowtex@tocattr{\the\attribute909}%
    \attribute909=\reflowtexLinkCount\relax
    \reflowtex@origcontentsline{#1}{#2}{#3}%
    \attribute909=\reflowtex@tocattr\relax
  \else
    \reflowtex@origcontentsline{#1}{#2}{#3}%
  \fi}
% the original, with the page number left out, and the fourth argument if
% the line had one (outside the conditional: the original reads its arguments)
\def\reflowtex@origcontentsline#1#2#3{%
  \def\reflowtex@tmpa{#3}%
  \ifx\reflowtex@tmpa\reflowtex@noarg \expandafter\@firstoftwo\else \expandafter\@secondoftwo\fi
  {\reflowtexOrigContentsline{#1}{#2}{}}{\reflowtexOrigContentsline{#1}{#2}{}{#3}}}
% With no page number, the dotted leaders of the standard classes would lead
% to nothing: spaced wider than any line, not one dot fits.
\AtBeginDocument{%
  \def\@dotsep{10000}%
  \let\reflowtexOrigContentsline\contentsline
  \let\contentsline\reflowtex@newcontentsline
  \let\reflowtex@starttoc\@starttoc
  \let\@starttoc\reflowtex@newstarttoc}
\makeatother

% The build pipeline compiles display-bearing snippets at several increasingly
% wide virtual measures. Paragraph node lists do not depend on this measure;
% finished display trees do, and comparing them lets the encoder recover their
% affine geometry without knowing whether the source used equation, align,
% gather, multline, or a package-defined display. This hook runs after the user
% preamble so it scales the final configured text width.
% The width increment the driver set (reflowtexWidthExtra, defined in
% reflowtex-capture.tex).
\AtBeginDocument{%
  % Wider for the display model's samples; narrower too, for the render
  % tests (tests/render), which compare the browser with TeX at other widths.
  \ifdim\reflowtexWidthExtra=0pt\else
    \advance\textwidth by \reflowtexWidthExtra
    \hsize=\textwidth
    \linewidth=\textwidth
    \columnwidth=\textwidth
  \fi
  \directlua{Serializer.note_source_width(\number\textwidth)}}

% TeX picks a display's skips from the line before it: \abovedisplayshortskip
% and \belowdisplayshortskip when that line ends left of the formula. The
% renderer re-breaks the line, so it has to make that choice itself. Every
% display is numbered here, its box stamped with the number, and what TeX
% decides with – \predisplaysize and the four skips – recorded under it
% (Serializer.note_display); the serializer puts them on the display item.
\newcount\reflowtexDisplayCount
\AtBeginDocument{\everydisplay\expandafter{\the\everydisplay
  \global\advance\reflowtexDisplayCount by 1
  \attribute912=\reflowtexDisplayCount
  \directlua{Serializer.note_display(\the\reflowtexDisplayCount)}%
  \aftergroup\reflowtexDisplayEnd}}
% When the display's math group closes: its leading is read no more.
\def\reflowtexDisplayEnd{\directlua{Serializer.note_display_end()}}
% LaTeX's \addvspace (a theorem's, list's or heading's space above) makes the
% space since the last box at least its argument, judged on \lastskip. After
% a display that skip is the below-display glue, which the renderer swaps for
% its own choice – so it must know the minimum too. Every \addvspace reports
% the \lastskip it found and what it asked for; the serializer attaches the
% minimum to the display whose skip it was (Serializer.note_addvspace). The
% argument may stretch (\addvspace{1em plus 1pt} above a ToC section entry):
% its natural width is what counts.
\let\reflowtexAddvspace\addvspace
\def\addvspace#1{%
  \ifvmode\directlua{Serializer.note_addvspace(\number\dimexpr\lastskip\relax, \number\glueexpr#1\relax)}\fi
  \reflowtexAddvspace{#1}}

% TikZ pictures, captured (the machinery is in reflowtex-capture.tex) when
% the preamble loaded TikZ.
\makeatletter
\@ifundefined{tikzpicture}{}{%
  \@ifundefined{tikzexternaldisable}{}{\tikzexternaldisable}%
  \let\reflowtexOrigTikzpicture\tikzpicture
  \let\reflowtexOrigEndtikzpicture\endtikzpicture
  \let\tikzpicture\reflowtex@tikzpicture
  \let\endtikzpicture\reflowtex@endtikzpicture}
\makeatother

% ── Cross-reference capture ──────────────────────────────────────────────────
% \ref and friends print a number; what the reader wants is to be taken to the
% thing numbered. Stamp every glyph of the printed text with a document-local
% link id and record which label that id points at, and leave an empty box where
% each \label stood so the target has a position in the flow.
%
% Nothing here changes what is typeset. Each reference still prints through the
% original macro, so its number, formatting and any package's decoration are
% unaffected, and the box carries no ink and no dimensions, so line breaking is
% unchanged.
%
% `"` is forced back to an ordinary character while this is read. A preamble may
% legitimately have made it active – the knowledge package's quotation syntax
% does, so that "term" is a reference – and an active quote inside the Lua string
% literals below is passed to Lua as a macro call rather than a delimiter, which
% fails at run time with a bare `unexpected symbol near '\'`. Catcodes are fixed
% when a definition is *read*, so the guard has to be here, not at the call.
\edef\reflowtexRestoreQuote{\catcode`\noexpand\"=\the\catcode`\"\relax}
\catcode`\"=12\relax
% The replacements are *defined* here and only swapped in inside the hook.
% Nothing with a parameter token may be written inside \AtBeginDocument: the
% hook system stores that code in a token list rather than a macro body, so `##`
% is not reduced to `#` there and the definitions would be read with doubled
% parameters.
\newcount\reflowtexLinkCount
\newcount\reflowtexAnchorCount
\makeatletter
% An empty, zero-sized box. A box rather than a kern or a whatsit so it survives
% as an ordinary list item the serializer can find by attribute, and zero-sized
% so it can neither show anything nor shift anything around it.
%
% Being zero-sized is not enough on a vertical list: TeX puts interline glue
% above every box there, from the depth of the line before, and the box then
% sets \prevdepth to its own depth, 0, so the next line's glue is computed as
% if nothing stood above it. A \label straight after \begin{theorem}, \item or
% \section would move every line after it. So in vertical mode the box goes in
% with no interline glue and \prevdepth put back: the next line is set exactly
% as without it.
\newbox\reflowtexAnchorBox
\def\reflowtex@anchorbox@{%
  \setbox\reflowtexAnchorBox=\hbox{}%
  \wd\reflowtexAnchorBox=0pt \ht\reflowtexAnchorBox=0pt \dp\reflowtexAnchorBox=0pt
  \box\reflowtexAnchorBox}
\def\reflowtex@anchorbox{%
  \ifvmode
    \begingroup
      \dimen@=\prevdepth \nointerlineskip
      \reflowtex@anchorbox@
      \prevdepth=\dimen@
    \endgroup
  \else
    \reflowtex@anchorbox@
  \fi}
% #1 is the saved original macro, #2 the label.
%
% One reference macro often prints through another – \eqref is (\ref{...}) – and
% the inner call would otherwise stamp the number with an id of its own, cutting
% the parentheses off from the digits they belong to. The outermost reference is
% the one the reader means, so an already-stamped group is left alone. An unset
% LuaTeX attribute reads as a large negative sentinel, hence the >0 test.
\def\reflowtex@wrapref#1#2{%
  \ifnum\attribute909>0
    #1{#2}%
  \else
    \global\advance\reflowtexLinkCount by 1
    \directlua{Serializer.note_link(\the\reflowtexLinkCount, "\luaescapestring{#2}")}%
    \begingroup\attribute909=\reflowtexLinkCount\relax #1{#2}\endgroup
  \fi}
\DeclareRobustCommand{\reflowtex@newref}[1]{\reflowtex@wrapref\reflowtex@ref{#1}}
\DeclareRobustCommand{\reflowtex@neweqref}[1]{\reflowtex@wrapref\reflowtex@eqref{#1}}
\DeclareRobustCommand{\reflowtex@newautoref}[1]{\reflowtex@wrapref\reflowtex@autoref{#1}}
\DeclareRobustCommand{\reflowtex@newnameref}[1]{\reflowtex@wrapref\reflowtex@nameref{#1}}
\DeclareRobustCommand{\reflowtex@newcref}[1]{\reflowtex@wrapref\reflowtex@cref{#1}}
\DeclareRobustCommand{\reflowtex@newCref}[1]{\reflowtex@wrapref\reflowtex@Cref{#1}}
% A destination the author wrote out in full. Hooked at hyperref's own
% \hyper@linkurl rather than at \href or \url: by the time hyperref calls it,
% the address has already been read with the catcodes a URL needs (% and # are
% ordinary there), so what arrives is the finished string. Re-reading the
% argument at \href would corrupt exactly the addresses that need the care.
\def\reflowtex@newlinkurl#1#2{%
  \ifnum\attribute909>0
    \reflowtex@linkurl{#1}{#2}%
  \else
    \global\advance\reflowtexLinkCount by 1
    \directlua{Serializer.note_link_url(\the\reflowtexLinkCount,
      "\luaescapestring{#2}")}%
    \begingroup\attribute909=\reflowtexLinkCount\relax
      \reflowtex@linkurl{#1}{#2}%
    \endgroup
  \fi}
% A destination with no scheme – \href{testmath/}{…}, \href{../book/#ch2}{…} –
% is a relative address on a web page, but hyperref takes it for a local file
% and turns it into a PDF path (testmath/.pdf). Hooked at \@hyper@linkfile,
% which still has the address as written (#4) and the anchor after its # (#3).
\def\reflowtex@newhyperlinkfile file:#1\\#2#3#4{%
  \ifnum\attribute909>0
    \reflowtex@hyperlinkfile file:#1\\{#2}{#3}{#4}%
  \else
    \global\advance\reflowtexLinkCount by 1
    \directlua{Serializer.note_link_file(\the\reflowtexLinkCount,
      "\luaescapestring{#4}", "\luaescapestring{#3}")}%
    \begingroup\attribute909=\reflowtexLinkCount\relax
      \reflowtex@hyperlinkfile file:#1\\{#2}{#3}{#4}%
    \endgroup
  \fi}
% Record and place the marker *before* delegating: packages that layer on \label
% – cleveref, knowledge – scan ahead from it for an optional argument, and
% anything left after the delegated call can be swallowed by that lookahead.
\DeclareRobustCommand{\reflowtex@newlabel@}[1]{%
  \global\advance\reflowtexAnchorCount by 1
  \directlua{Serializer.note_label(\the\reflowtexAnchorCount, "\luaescapestring{#1}")}%
  \begingroup\attribute910=\reflowtexAnchorCount\relax
    \reflowtex@anchorbox
  \endgroup
  \reflowtex@label{#1}}
% Installed at \begin{document}, and this hook is declared after the user
% preamble so it runs after theirs: hyperref redefines \ref itself at
% \begin{document}, and a preamble-time redefinition is simply overwritten.
% The label is captured here, at the call site, rather than taken from
% hyperref's anchor: under `hypertexnames=false` the anchor is an opaque
% counter, while the label is what a publishing toolchain can still match
% against the labels other chapters define.
\AtBeginDocument{%
  \let\reflowtex@ref\ref     \let\ref\reflowtex@newref
  \@ifundefined{eqref}{}{\let\reflowtex@eqref\eqref \let\eqref\reflowtex@neweqref}%
  \@ifundefined{autoref}{}{\let\reflowtex@autoref\autoref
                           \let\autoref\reflowtex@newautoref}%
  \@ifundefined{nameref}{}{\let\reflowtex@nameref\nameref
                           \let\nameref\reflowtex@newnameref}%
  % cleveref's \cref/\Cref take a comma-separated list; the whole printed run
  % becomes one link to the first label named, which is where it would jump.
  \@ifundefined{cref}{}{\let\reflowtex@cref\cref \let\cref\reflowtex@newcref}%
  \@ifundefined{Cref}{}{\let\reflowtex@Cref\Cref \let\Cref\reflowtex@newCref}%
  \@ifundefined{hyper@linkurl}{}{\let\reflowtex@linkurl\hyper@linkurl
                                 \let\hyper@linkurl\reflowtex@newlinkurl}%
  \@ifundefined{@hyper@linkfile}{}{\let\reflowtex@hyperlinkfile\@hyper@linkfile
                                   \let\@hyper@linkfile\reflowtex@newhyperlinkfile}%
  \let\reflowtex@label\label \let\label\reflowtex@newlabel@
  % A \label inside a display never reaches \label: amsmath swaps in
  % \label@in@display, which only stores the name and replays it later through
  % \ltx@label – a copy of \label taken when amsmath was loaded, long before
  % this hook runs. Redirect that copy as well, or every equation label is lost.
  \@ifundefined{ltx@label}{}{\let\ltx@label\reflowtex@newlabel@}}
\makeatother
\reflowtexRestoreQuote

% ── Outline capture ──────────────────────────────────────────────────────────
% The document's outline – its sections and subsections, and its theorems,
% lemmas and the like – for a page to show as a table of contents. Each entry
% is recorded (Serializer.note_outline) with an anchor at its place in the
% flow, labelled reflowtex-outline-N, which becomes an element the page can
% scroll to. Nothing here changes what is typeset.
%
% Sections: whatever goes into the table of contents, caught at
% \addcontentsline{toc}, whose text already carries the number
% (\numberline{3.1}Title). Its anchor goes into the vertical list after the
% heading, with no interline glue and \prevdepth restored, so the spacing to
% the next line is untouched.
%
% Unnumbered sections (\section*, and so the bibliography's References):
% they never reach the table of contents, so they are caught at \@ssect,
% with the level remembered from \@startsection. Their anchor goes in at the
% start of the heading's own paragraph (a one-shot para/begin hook). A
% starred section a document also adds to the contents by hand is listed
% once (the serializer drops the repeat).
%
% Theorem-like environments: caught at \@begintheorem (amsthm's, or the
% kernel's pair), after the heading is set up. The anchor goes in at the
% start of the theorem's first paragraph (a one-shot para/begin hook), so a
% jump lands on the theorem itself – without starting that paragraph any
% earlier than LaTeX would: a \label right after \begin{thm} stays in
% vertical mode, and the line end after it stays a non-space.
\newcount\reflowtexOutlineCount
\makeatletter
% #1 kind, #2 level, #3 number, #4 title (macros holding the text). Records
% the entry; its anchor is \reflowtexAnchorCount afterwards.
\def\reflowtex@outline@note#1#2#3#4{%
  \global\advance\reflowtexAnchorCount by 1
  \global\advance\reflowtexOutlineCount by 1
  \directlua{Serializer.note_label(\the\reflowtexAnchorCount,
      "reflowtex-outline-\the\reflowtexOutlineCount")
    Serializer.note_outline(\the\reflowtexAnchorCount,
      "\luaescapestring{\detokenize\expandafter{#1}}", #2,
      "\luaescapestring{\detokenize\expandafter{#3}}",
      "\luaescapestring{\detokenize\expandafter{#4}}")}}
% Records the entry and places its anchor here.
\def\reflowtex@outline@mark#1#2#3#4{%
  \reflowtex@outline@note{#1}{#2}{#3}{#4}%
  \begingroup\attribute910=\reflowtexAnchorCount\relax\reflowtex@anchorbox\endgroup}
\def\reflowtex@outline@acl#1#2#3{%
  \def\reflowtex@ol@file{#1}\def\reflowtex@ol@toc{toc}%
  \ifx\reflowtex@ol@file\reflowtex@ol@toc
    \begingroup
      \protected@edef\reflowtex@ol@kind{#2}%
      \protected@edef\reflowtex@ol@text{#3}%
      \let\reflowtex@ol@none\@empty
      \reflowtex@outline@mark\reflowtex@ol@kind{0}\reflowtex@ol@none\reflowtex@ol@text
    \endgroup
    \reflowtex@toc@anchor
  \fi
  \reflowtex@acl{#1}{#2}{#3}}
% The anchor a table of contents' k-th line links to (see \contentsline above).
\def\reflowtex@toc@anchor{%
  \global\advance\reflowtexTocAnchorCount by 1
  \global\advance\reflowtexAnchorCount by 1
  \directlua{Serializer.note_label(\the\reflowtexAnchorCount,
    "reflowtex-toc-\the\reflowtexTocAnchorCount")}%
  \begingroup\attribute910=\reflowtexAnchorCount\relax\reflowtex@anchorbox\endgroup}
% #1 heading name (Theorem), #2 number, #3 note (may be empty).
\def\reflowtex@outline@thm#1#2#3{%
  \begingroup
    \protected@edef\reflowtex@ol@name{#1}%
    \protected@edef\reflowtex@ol@num{#2}%
    \protected@edef\reflowtex@ol@note{#3}%
    \protected@edef\reflowtex@ol@title{\reflowtex@ol@name
      \ifx\reflowtex@ol@note\@empty\else\space(\reflowtex@ol@note)\fi}%
    \protected@edef\reflowtex@ol@kind{\@currenvir}%
    \reflowtex@outline@note\reflowtex@ol@kind{9}\reflowtex@ol@num\reflowtex@ol@title
    \edef\reflowtex@ol@hook{\noexpand\AddToHookNext{para/begin}{%
      \begingroup\attribute910=\the\reflowtexAnchorCount\relax
      \noexpand\reflowtex@anchorbox\endgroup}}%
  \expandafter\endgroup\reflowtex@ol@hook}
% #1 the heading's text (unnumbered section).
\def\reflowtex@outline@starred#1{%
  \begingroup
    \protected@edef\reflowtex@ol@title{#1}%
    \let\reflowtex@ol@none\@empty
    \reflowtex@outline@note\reflowtex@ol@secname{0}\reflowtex@ol@none\reflowtex@ol@title
    \edef\reflowtex@ol@hook{\noexpand\AddToHookNext{para/begin}{%
      \begingroup\attribute910=\the\reflowtexAnchorCount\relax
      \noexpand\reflowtex@anchorbox\endgroup}}%
  \expandafter\endgroup\reflowtex@ol@hook}
\let\reflowtex@ol@secname\@empty
% The wrappers are defined here and only swapped in at \begin{document} (a
% hook's code cannot hold parameter tokens; see the cross-reference section).
\def\reflowtex@ol@startsection#1#2#3#4#5#6{%
  \def\reflowtex@ol@secname{#1}\reflowtex@ss{#1}{#2}{#3}{#4}{#5}{#6}}
\def\reflowtex@ol@ssect#1#2#3#4#5{%
  \reflowtex@outline@starred{#5}\reflowtex@sss{#1}{#2}{#3}{#4}{#5}}
% amsthm's head reads the note as a delimited argument, [#3], which takes off
% the braces an author puts round a note holding a `]` (\cite[p.~3]{key});
% passed on as [#3], the head would end the note at that `]`. So it is passed
% on braced, [{#3}], and the head takes the braces off again. (LaTeX's own
% theorems pass the note as {#3}, and need nothing.)
\def\reflowtex@ams@bt#1#2[#3]{\reflowtex@bt{#1}{#2}[{#3}]\reflowtex@outline@thm{#1}{#2}{#3}}
\def\reflowtex@k@bt#1#2{\reflowtex@bt{#1}{#2}\reflowtex@outline@thm{#1}{#2}{}}
\def\reflowtex@k@obt#1#2#3{\reflowtex@obt{#1}{#2}{#3}\reflowtex@outline@thm{#1}{#2}{#3}}
\AtBeginDocument{%
  \let\reflowtex@acl\addcontentsline
  \let\addcontentsline\reflowtex@outline@acl
  \let\reflowtex@ss\@startsection  \let\@startsection\reflowtex@ol@startsection
  \let\reflowtex@sss\@ssect        \let\@ssect\reflowtex@ol@ssect
  \let\reflowtex@bt\@begintheorem
  \@ifpackageloaded{amsthm}%
    {\let\@begintheorem\reflowtex@ams@bt}%
    {\let\@begintheorem\reflowtex@k@bt
     \let\reflowtex@obt\@opargbegintheorem
     \let\@opargbegintheorem\reflowtex@k@obt}}
\makeatother

% ── Footnote capture ─────────────────────────────────────────────────────────
% Install this after the user preamble: packages such as hyperref legitimately
% replace \@footnotemark, and wrapping an earlier definition would be lost.
% TeX places the marker in the paragraph but sends its body to an insertion.
% Give both halves one document-local id; separate attributes keep the body's
% own leading number from recursively acting as another tooltip source.
\makeatletter
\newcount\reflowtexFootnoteCount
\newcount\reflowtexCurrentFootnote
% At the end of \begin{document}: with tagging on, LaTeX's footnote code
% (latex-lab-footnotes) sets \@footnotemark and \@footnotetext as the
% document begins, and wrapping the definitions before would be lost.
\AddToHook{begindocument/end}{%
  \let\reflowtexOrigFootnotemark\@footnotemark
  % (the attribute set and unset, no group: with tagging, LaTeX finishes
  % the mark later, and needs what it began with – \@x@sf – then)
  \def\@footnotemark{%
    \global\advance\reflowtexFootnoteCount by 1
    \global\reflowtexCurrentFootnote=\reflowtexFootnoteCount
    \attribute906=\reflowtexCurrentFootnote\relax
    \reflowtexOrigFootnotemark
    \attribute906=-"7FFFFFFF\relax}%
  \let\reflowtexOrigFootnotetext\@footnotetext
  \long\def\@footnotetext#1{%
    \attribute907=\reflowtexCurrentFootnote\relax
    \reflowtexOrigFootnotetext{#1}%
    \attribute907=-"7FFFFFFF\relax}}
\makeatother

% ── Margin notes ─────────────────────────────────────────────────────────────
% \marginpar[left]{right} as an aside of kind "marginpar", place=margin: the
% viewer sets it in the margin by its line, or behind a mark where there is
% none. There is one margin on the web, so the text for a left margin is not
% used. After the user preamble, so a class's own \marginpar is replaced too.
\renewcommand\marginpar[2][]{\reflowtexAside{marginpar}{place=margin}{#2}}
